% GENERATED FILE. Edit canonical lessons, then run npm run generate:books.
% canonical-source-hash: fnv1a64:dead0a6ea3620f0a
% canonical-lessons: GE-C158-angebot, GE-C158-rechnung, GE-C158-tuete, GE-C158-paket, GE-C158-briefmarke

\chapter{An offer, A bill, A bag, A parcel, A stamp}
\label{ch:ge-an-offer-a-bill-a-bag-a-parcel-a-stamp}
\hlchaptermodality{158}

\begin{chapteropening}
\textbf{By the end of this chapter:} \emph{I can say an offer, a bill, a bag, a parcel and a stamp.}

Complete the last lesson of chapter 158: I can say an offer, a bill, a bag, a parcel and a stamp.
\end{chapteropening}

\section[das Angebot]{das Angebot — an offer}

\label{lesson:GE-C158-angebot}



\begin{quote}
Before the new one: say the German for a credit card, then the German for a price.
\end{quote}

\subsection*{You'll want to know: das Angebot}
\textbf{das Angebot} — \textquotedblleft{}an offer\textquotedblright{}.

\begin{cousinweb}
\textbf{Angebot} is the noun of the verb \textquotedblleft{}to offer\textquotedblright{}: \emph{an-} + a verb meaning \textquotedblleft{}to offer, to bid\textquotedblright{}, which EtymDB 2.1 lists beside English \emph{bid} and Dutch \emph{bieden}.
\end{cousinweb}

\begin{grammarlens}[title={What you've built}]
One more word for this chapter.
\end{grammarlens}

\subsection*{Your turn}
\begin{itemize}
\raggedright
  \item \emph{Say it:} \emph{das Angebot}
  \item \emph{Say it:} \emph{das Angebot}, once more
  \item \emph{Say it:} say \emph{das Angebot}
  \item \emph{Recall:} say the German for a credit card, then the German for a price, then say \emph{das Angebot} again
\end{itemize}

\begin{tcolorbox}[breakable,colback=teal!4,colframe=teal!35!black,title={Before you move on}]
What does das Angebot mean? (\textquotedblleft{}An offer\textquotedblright{}.) Say it once more.
\end{tcolorbox}
\section[die Rechnung]{die Rechnung — a bill}

\label{lesson:GE-C158-rechnung}



\begin{quote}
Before the new one: say the German for a price, then the German for an offer.
\end{quote}

\subsection*{You'll want to know: die Rechnung}
\textbf{die Rechnung} — \textquotedblleft{}a bill\textquotedblright{}.

\begin{cousinweb}
\textbf{Rechnung} is the verb \textquotedblleft{}to reckon\textquotedblright{}, Old High German \emph{rehhanōn} $\leftarrow$ Proto-Germanic \emph{*rekanōną}, + \emph{-ung} — EtymDB 2.1 lists it beside English \emph{reckoning}.
\end{cousinweb}

\begin{grammarlens}[title={What you've built}]
One more word for this chapter.
\end{grammarlens}

\subsection*{Your turn}
\begin{itemize}
\raggedright
  \item \emph{Say it:} \emph{die Rechnung}
  \item \emph{Say it:} \emph{die Rechnung}, once more
  \item \emph{Say it:} say \emph{die Rechnung}
  \item \emph{Recall:} say the German for a price, then the German for an offer, then say \emph{die Rechnung} again
\end{itemize}

\begin{tcolorbox}[breakable,colback=teal!4,colframe=teal!35!black,title={Before you move on}]
What does die Rechnung mean? (\textquotedblleft{}A bill\textquotedblright{}.) Say it once more.
\end{tcolorbox}
\section[die Tüte]{die Tüte — a bag (paper or plastic)}

\label{lesson:GE-C158-tuete}



\begin{quote}
Before the new one: say the German for an offer, then the German for a bill.
\end{quote}

\subsection*{You'll want to know: die Tüte}
\textbf{die Tüte} — \textquotedblleft{}a bag (paper or plastic)\textquotedblright{}.

\begin{cousinweb}
\textbf{Tüte} $\leftarrow$ Middle Low German \emph{tūte} (EtymDB 2.1) — borrowed from Low German; EtymDB lists Dutch \emph{tuit}, \textquotedblleft{}spout\textquotedblright{}, beside it. A paper bag was rolled into a cone.
\end{cousinweb}

\begin{grammarlens}[title={What you've built}]
One more word for this chapter.
\end{grammarlens}

\subsection*{Your turn}
\begin{itemize}
\raggedright
  \item \emph{Say it:} \emph{die Tüte}
  \item \emph{Say it:} \emph{die Tüte}, once more
  \item \emph{Say it:} say \emph{die Tüte}
  \item \emph{Recall:} say the German for an offer, then the German for a bill, then say \emph{die Tüte} again
\end{itemize}

\begin{tcolorbox}[breakable,colback=teal!4,colframe=teal!35!black,title={Before you move on}]
What does die Tüte mean? (\textquotedblleft{}A bag (paper or plastic)\textquotedblright{}.) Say it once more.
\end{tcolorbox}
\section[das Paket]{das Paket — a parcel}

\label{lesson:GE-C158-paket}



\begin{quote}
Before the new one: say the German for a bill, then the German for a bag.
\end{quote}

\subsection*{You'll want to know: das Paket}
\textbf{das Paket} — \textquotedblleft{}a parcel\textquotedblright{}.

\begin{cousinweb}
\textbf{Paket} $\leftarrow$ French \emph{paquet}, \textquotedblleft{}packet\textquotedblright{}, from Dutch \emph{pak}, \textquotedblleft{}bundle\textquotedblright{} (EtymDB 2.1) — borrowed. English \emph{packet} is the same French word.
\end{cousinweb}

\begin{grammarlens}[title={What you've built}]
One more word for this chapter.
\end{grammarlens}

\subsection*{Your turn}
\begin{itemize}
\raggedright
  \item \emph{Say it:} \emph{das Paket}
  \item \emph{Say it:} \emph{das Paket}, once more
  \item \emph{Say it:} say \emph{das Paket}
  \item \emph{Recall:} say the German for a bill, then the German for a bag, then say \emph{das Paket} again
\end{itemize}

\begin{tcolorbox}[breakable,colback=teal!4,colframe=teal!35!black,title={Before you move on}]
What does das Paket mean? (\textquotedblleft{}A parcel\textquotedblright{}.) Say it once more.
\end{tcolorbox}
\section[die Briefmarke]{die Briefmarke — a stamp}

\label{lesson:GE-C158-briefmarke}



\begin{quote}
Before the new one: say the German for a bag, then the German for a parcel.
\end{quote}

\subsection*{You'll want to know: die Briefmarke}
\textbf{die Briefmarke} — \textquotedblleft{}a stamp\textquotedblright{}.

\begin{cousinweb}
\textbf{Briefmarke} is the word for \textquotedblleft{}letter\textquotedblright{} + \emph{Marke}, \textquotedblleft{}token, mark\textquotedblright{} $\leftarrow$ French \emph{marque} (EtymDB 2.1): a token stuck on a letter.
\end{cousinweb}

\begin{grammarlens}[title={What you've built}]
One more word for this chapter.
\end{grammarlens}

\subsection*{Your turn}
\begin{itemize}
\raggedright
  \item \emph{Say it:} \emph{die Briefmarke}
  \item \emph{Say it:} \emph{die Briefmarke}, once more
  \item \emph{Say it:} say \emph{die Briefmarke}
  \item \emph{Recall:} say the German for a bag, then the German for a parcel, then say \emph{die Briefmarke} again
\end{itemize}

\begin{tcolorbox}[breakable,colback=teal!4,colframe=teal!35!black,title={Before you move on}]
What does die Briefmarke mean? (\textquotedblleft{}A stamp\textquotedblright{}.) Say it once more.
\end{tcolorbox}
